\documentclass[11pt]{article}
%\usepackage{../../styles/activity}
\usepackage{c://pctex/beginact}

\usepackage{xr}
\externaldocument{0-MR}

\lhead{}
%\chead{\textbf{\Large{\hspace{0pt}Beginning Activities for Section~2.2}}\\\beghead}
\bahead{2.2}
\rhead{}
\lfoot{}
\rfoot{}
\cfoot{\hspace{0pt}\scalebox{0.4}{\includegraphics{cc-by-nc-sa.eps}}}

\begin{document}

\subsection*{Beginning Activity 1 (Logically Equivalent Statements)}
\begin{enumerate}

\item 
\begin{tabular}[t]{| c | c || c | c | c | c | c |}  \hline
$P$  &  $Q$  &  $P \wedge Q$  &  $\mynot  ( {P \wedge Q} )$  &  $\mynot P$  &  
$\mynot Q$  &  $\mynot  P \vee \mynot  Q$ \\ \hline
T  &  T  &  T  &  F  &  F  &  F  &  F  \\ \hline
T  &  F  &  F  &  T  &  F  &  T  &  T  \\ \hline
F  &  T  &  F  &  T  &  T  &  F  &  T  \\ \hline
F  &  F  &  F  &  T  &  T  &  T  &  T  \\ \hline
\end{tabular}

\item The truth table shows that the expressions $\mynot  ( {P \wedge Q} )$ and 
$\mynot  P \vee \mynot  Q$ are logically equivalent.

\item Using the fact that $\mynot  ( {P \wedge Q} )$ and $\mynot  P \vee \mynot  Q$ are logically equivalent, we see that the negation of ``I will play golf and I will mow the lawn'' is ``I will not play golf or I will not mow the lawn.''

\item Different definitions for  $P$ and  $Q$  other than the ones shown here can be used.
Let  $P$  be the statement ``You do not clean your room.''  Let  $Q$  be the statement  ``You cannot watch TV.''
The first statement is then $P \to Q$  and the second statement is  $\mynot P \vee Q$. 
\item 
\begin{tabular}[t]{| c | c || c | c | c |}  \hline
$P$  &  $Q$  &  $\mynot P$  &  $P \to Q$  &  $\mynot P \vee Q$ \\ \hline
T  &  T  &  F  &  T  &  T  \\ \hline
T  &  F  &  F  &  F  &  F  \\ \hline
F  &  T  &  T  &  T  &  T  \\ \hline
F  &  F  &  T  &  T  &  T  \\ \hline
\end{tabular}

Therefore, $P \to Q$ is logically equivalent to  $\mynot  P \vee Q$, or symbolically,  
$( {P \to Q} ) \equiv ( {\mynot  P \vee Q} )$.

\item Statement 1 and Statement 2 are logically equivalent.  The conditional statement is false when $P$ is true and $Q$ is false.  So when the conditional statement is false, the statement ``You do not clean your room and you can watch TV'' is true.  Symbolically, this statement is $P \wedge \mynot Q$ . 

\item 
\begin{tabular}[t]{| c | c || c | c | c | c |}  \hline
$P$  &  $Q$  &  $P \to Q$  &  $\mynot ( {P \to Q} )$  &  $\mynot Q$ &  $P \wedge \mynot Q$ \\ \hline
T  &  T  &  T  &  F  &  F & F \\ \hline
T  &  F  &  F  &  T  &  T & T \\ \hline
F  &  T  &  T  &  F  &  F & F  \\ \hline
F  &  F  &  T  &  F  &  T & F \\ \hline
\end{tabular}

This shows that  $\mynot  ( {P \to Q} )$  is logically equivalent to  
$P \wedge \mynot  Q$.

\end{enumerate}
\hbreak


\newpage
\subsection*{ Beginning Activity 2 (Converse and Contrapositive)}

\begin{enumerate}
\item Statements (a) and (c) are true.  Statements (b) and (d) are false.

\item The converse of ``If $x=3$, then $x^2 = 9$'' is ``If $x^2=9$, then $x=3$'', which is Statement (b).

The contrapositive of ``If $x=3$, then $x^2 = 9$'' is ``If  $x^2 \ne 9$, then $x \ne 3$'', which is Statement (c).

\item
\begin{tabular}[t]{| c | c || c | c | c | c | c |}  \hline
$P$  &  $Q$  &  $P \to Q$  &  $Q \to P$  &  $\mynot Q$  &  $\mynot P$  &  
$\mynot  Q \to\mynot  P$ \\ \hline
T  &  T  &  T  &  T  &  F  &  F  &  T  \\ \hline
T  &  F  &  F  &  T  &  T  &  F  &  F  \\ \hline
F  &  T  &  T  &  F  &  F  &  T  &  T  \\ \hline
F  &  F  &  T  &  T  &  T  &  T  &  T  \\ \hline
\end{tabular}

The columns for $P \to Q$ and $\mynot  Q \to\mynot  P$ show that these two statements are logically equivalent.  The columns for $P \to Q$ and $Q \to P$ show that these two statements are not logically equivalent.

It is important to be aware that the converse of a conditional statement is not logically equivalent to the conditional statement.  However, the contrapositive of the conditional statement is logically equivalent to the conditional statement.  We will use this to help use prove certain types of conditional statements in Section~3.2.
\end{enumerate}
\hbreak



\end{document}

